\chapter{Phase 2 --- Microarchitecture and RTL Design}\label{ch:rtl}

\weeks{6--12} \quad Milestone M2 (together with \cref{ch:verification})

In this phase the specification becomes hardware. You first design the
\emph{microarchitecture}: which blocks exist, how data moves between them, and
what happens on each clock cycle of a frame. Then you write it in synthesizable
SystemVerilog, one block at a time, each verified against the golden model
before you move on (\cref{ch:verification}). Build and test go together: no
block is \enquote{done} until its test passes.

\begin{whybox}
Phase~2 is where you learn the core skill of VLSI DSP: mapping an algorithm onto
time and hardware resources. It produces the cycle counts that test H4 and the
first resource numbers for RQ3. Most of the interesting decisions already have
a cost model, from \cref{ch:workload}, and a specification, from
\cref{ch:codesign}. Here you make them concrete.
\end{whybox}

\begin{center}\small
\begin{tabularx}{\textwidth}{@{}LL@{}}
\toprule
\textbf{Inputs} & \textbf{Outputs} \\
\midrule
\code{spec-v1}: formats, tables, frame configuration; golden vectors; the
cocotb setup from Phase~0
& Lint-clean RTL in \file{rtl/}; every block passing its tests; a resource log
(Yosys and Vivado out-of-context synthesis) per milestone commit; a block
diagram and frame schedule for the report \\
\bottomrule
\end{tabularx}
\end{center}

% =========================================================================
\section{Top-level architecture}\label{sec:rtl:arch}

\Cref{fig:rtl:arch} shows the datapath. Audio samples stream in at the sample
rate. Frames are processed in bursts on a much faster system clock. Parameters
arrive over a bus from the control plane.

\begin{figure}[htbp]
\centering
\begin{tikzpicture}[scale=0.8,transform shape,font=\sffamily\scriptsize,
  b/.style={draw=brand,fill=brandlight,rounded corners=1mm,minimum height=8mm,minimum width=16mm,align=center},
  m/.style={draw=conceptc,fill=conceptc!8,rounded corners=1mm,minimum height=7mm,minimum width=13mm,align=center},
  c/.style={draw=rqc,fill=rqc!6,rounded corners=1mm,minimum height=7mm,minimum width=16mm,align=center},
  a/.style={-{Stealth[length=1.8mm]},line width=0.6pt},
  ca/.style={-{Stealth[length=1.8mm]},line width=0.6pt,draw=rqc,dashed}]
  % data plane
  \node[b] (in)  at (0,0)    {\texttt{stft\_frontend}\\window, framing};
  \node[b] (fft) at (2.4,0)  {\texttt{fft\_r2\_mem}\\FFT and IFFT};
  \node[b,minimum width=86mm,minimum height=11mm] (bp) at (8.2,0) {\texttt{bin\_pipeline} --- IMCRA + parametric Wiener, one bin per cycle};
  \node[b] (ola) at (2.4,-1.9) {\texttt{ola\_backend}\\window, overlap-add};
  \draw[a] (-1.7,0) node[left]{samples in} -- (in);
  \draw[a] (in) -- (fft);
  \draw[a] ([yshift=1.5mm]fft.east) -- node[above,font=\sffamily\tiny]{$X_k$} ([yshift=1.5mm]bp.west);
  \draw[a] ([yshift=-1.5mm]bp.west) -- node[below,font=\sffamily\tiny]{$G_kX_k$} ([yshift=-1.5mm]fft.east);
  \draw[a] (fft) -- node[right,font=\sffamily\tiny]{IFFT output} (ola);
  \draw[a] (ola.south) |- (-1.7,-2.95) node[left]{samples out};
  % helpers of the per-bin pipeline
  \foreach \n/\x/\t in {st/4.7/{\texttt{state\_ram}\\23 words/bin},
                        mt/6.45/{\texttt{min\_tracker}\\$\times2$},
                        rc/8.2/{\texttt{recip\_unit}\\\texttt{exp\_unit}},
                        gl/9.95/{\texttt{gain\_lut}\\2 banks},
                        tg/11.7/{\texttt{tilt\_gain}\\$g_\text{out}r^k$}}{
    \node[m] (\n) at (\x,-1.9) {\t};
    \draw[conceptc,line width=0.5pt] (\n.north) -- (\n.north |- bp.south);
  }
  % control
  \node[c] (fc) at (0,-4.0)   {\texttt{frame\_ctrl}\\schedule, counters};
  \node[c] (rg) at (9.95,-4.0) {\texttt{regs\_axil}\\AXI4-Lite};
  \draw[ca] (rg.north) -- (gl.south);
  \draw[ca] (rg.east) -| (tg.south);
  \draw[ca] (rg.west) -- node[below]{CTRL, STATUS, cycle counters} (fc.east);
  \draw[a,draw=rqc] (13.4,-4.0) node[right,text=rqc,align=left]{control\\plane} -- (rg.east);
\end{tikzpicture}
\caption{Microarchitecture of \texttt{ise\_top}. Blue: data plane. Grey: helper
units and memories of the per-bin pipeline. Purple, dashed: control. One FFT
unit is shared between the forward and inverse transforms.}
\label{fig:rtl:arch}
\end{figure}

Three decisions shape this architecture:
\begin{itemize}
  \item \textbf{One clock domain.} The whole datapath runs on one system clock.
  A \emph{sample strobe}, a one-cycle pulse at the sample rate generated by
  \code{frame\_ctrl}, tells the front end when a new sample is valid. Crossing to
  the audio codec's own clock is handled at the edge of the FPGA design
  (\cref{ch:fpga}), not inside \code{ise\_top}.
  \item \textbf{One shared FFT.} Forward and inverse transforms never overlap
  within a frame, so one \code{fft\_r2\_mem} does both. The inverse uses the
  conjugation trick from \cref{ch:codesign}.
  \item \textbf{Bins stream in order.} The per-bin pipeline consumes one bin per
  cycle, from $k=0$ to $K-1$. Almost everything is independent across bins,
  and the 5-tap frequency smoothing becomes a sliding window over a short shift
  register.
\end{itemize}

% =========================================================================
\section{The frame schedule}\label{sec:rtl:schedule}

A frame is processed in a burst of cycles, while the next block of samples is
captured into the other half of a \emph{ping-pong} buffer
(\cref{fig:rtl:schedule}). The estimate below assumes the thesis frame
configuration; for the configuration you chose in Phase~1, scale it by the
number of bins.

\begin{figure}[htbp]
\centering
\begin{tikzpicture}[font=\sffamily\scriptsize,x=0.86mm,y=6mm]
  \foreach \s/\e/\l/\c in {0/10/{window, load}/brand, 10/68/{FFT ($\approx$5.6k)}/stepc, 68/80/{bins ($\approx$1.05k)}/whyc, 80/138/{IFFT ($\approx$5.6k)}/stepc, 138/150/{OLA}/brand}{
    \fill[\c!55,rounded corners=0.5mm] (\s,0) rectangle (\e,0.8);
    \node[font=\sffamily\tiny] at ({(\s+\e)/2},0.4) {\l};
  }
  \draw[-{Stealth}] (0,-0.2) -- (160,-0.2) node[right]{cycles};
  \node[anchor=north] at (0,-0.25) {0};
  \node[anchor=north] at (150,-0.25) {$\approx$13.5k};
  \fill[conceptc!25] (0,1.2) rectangle (160,1.8);
  \node at (80,1.5) {meanwhile: capture the next 512 samples into the other buffer half (one sample per strobe)};
\end{tikzpicture}
\caption{One frame of processing (estimate, thesis configuration, one butterfly
per cycle, one bin per cycle). At \qty{5}{\MHz} the budget is \num{53333} cycles
per frame, so the datapath is busy for about a quarter of the frame and can be
clock-gated for the rest.}
\label{fig:rtl:schedule}
\end{figure}

This is the calculation behind hypothesis H4. Make \code{frame\_ctrl} count the
cycles of each phase in hardware and expose the total in a debug register. That
makes H4 a measurement rather than an estimate.

% =========================================================================
\section{Designing the blocks}\label{sec:rtl:blocks}

\subsection{The FFT (\texttt{fft\_r2\_mem})}

A \emph{memory-based} radix-2 FFT keeps all $N$ complex values in a dual-port
RAM and runs one butterfly per cycle for $\log_2 N$ stages. It is the smallest
FFT architecture and is fast enough here: \num{5120} butterflies per 1024-point
transform. Pipelined architectures such as R2SDF or
R2\textsuperscript{2}SDF~\cite{he1996r22sdf,ch12-13:garrido2022fft} are faster,
but they are designed for continuous streaming, which you do not need.

You need four pieces:
\begin{itemize}
  \item \textbf{Address generation.} In stage $s$, butterfly $j$ reads two
  addresses that differ in bit $s$. A counter plus a bit-insertion gives both.
  \item \textbf{Twiddle ROM.} $W_N^m=e^{-j2\pi m/N}$. Store a quarter cosine
  wave and derive sine and cosine for all $m$ from symmetry.
  \item \textbf{Bit-reversed order.} Either write the input in bit-reversed
  order or read the output that way. Reversing an address is just wiring.
  \item \textbf{Block floating point.} Each stage can grow values by one bit.
  Detect growth into the guard bits during stage $s$, apply a right shift to the
  whole block in stage $s+1$, and count the shifts in a shared exponent. This
  must follow \file{golden/fx_stft.py} exactly.
\end{itemize}

Start with the simplest correct version: a complex FFT that takes real input
with zero imaginary part. It wastes half its work, but it is easy to verify.
Only once it is bit-exact, implement the real-FFT packing (an $N$-point real FFT
computed as an $N/2$-point complex FFT plus a post-processing pass), and only if
Phase~1 modelled it.

\subsection{The per-bin pipeline (\texttt{bin\_pipeline})}

Map the Python code line by line onto pipeline stages (\cref{tab:rtl:stages}).
Each stage holds only a few operations, so even long operations like the
reciprocal can be pipelined to accept one bin per cycle.

\begin{table}[htbp]
\centering\footnotesize
\caption{From Python lines to pipeline stages (an example grouping; your
spec decides the details).}
\label{tab:rtl:stages}
\begin{tabularx}{\textwidth}{@{}lLl@{}}
\toprule
Stage & Work & Thesis source \\
\midrule
1 & Read state word for bin $k$; compute $|X_k|^2$ (or its log); push into the 5-bin window & \file{pipeline.py:179} \\
2 & Frequency smoothing over bins $k-2\ldots k+2$; recursive smoothing of $S$ & \file{imcra.py:117-118} \\
3 & Minimum tracker 1; $\gamma_\text{min}$, $\zeta$ (log domain: subtractions); indicator $I_k$ & \file{imcra.py:119-126} \\
4 & Second smoothing over $I\cdot P$ and $I$ for bins $k-2\ldots k+2$; one reciprocal & \file{imcra.py:129-133} \\
5 & Minimum tracker 2; $\tilde\gamma_\text{min}$, $\tilde\zeta$; $q$ & \file{imcra.py:134-146} \\
6 & $\gamma$, $\nu$ (cached $H$, O3); $p$ via \code{exp\_unit} or the sigmoid rewrite; $\lambda_d$ update & \file{imcra.py:149-164} \\
7 & Wiener: decision-directed $\xi$ (shared reciprocal of the noise PSD) & \file{wiener.py:89-98} \\
8 & Gain: \code{gain\_lut}[index($\xi$)], times tilt and output gain from \code{tilt\_gain}, clip & \file{wiener.py:102-112} \\
9 & $G_kX_k$ to the IFFT input; previous clean power; write state back & \file{wiener.py:115} \\
\bottomrule
\end{tabularx}
\end{table}

Two details need care:
\begin{itemize}
  \item \textbf{Look-ahead.} Smoothing at bin $k$ needs bins $k+1$ and $k+2$. The
  first iteration therefore runs two bins \emph{behind} the input, and the second
  smoothing, which needs the indicator of bins up to $k+2$, two bins behind
  that. Implement both with short shift registers. At the spectrum edges,
  \code{np.convolve(\ldots, mode="same")} treats missing bins as zero, so your
  hardware must do exactly the same.
  \item \textbf{State memory.} Each bin's state is read once and written once per
  frame, at a different address from every other bin. There are therefore no
  read-after-write hazards between bins, only a latency to cover. Store each
  bin's state as one wide word, or as a few parallel memories, so the read takes
  one cycle.
\end{itemize}

\subsection{Arithmetic and table units}

\begin{description}[style=nextline,leftmargin=1.2em]
  \item[\texttt{recip\_unit}] Normalise the input with a leading-zero count, look
  up a seed for $1/d$ from the top mantissa bits, then do one Newton--Raphson
  step $y_1=y_0(2-dy_0)$, which costs two multiplies. Pipeline it so it accepts
  one input per cycle. Share it between the denominators that occur in the same
  stage.
  \item[\texttt{exp\_unit}] A table over $\nu\in[0,16]$, with or without linear
  interpolation, as specified.
  \item[\texttt{gain\_lut}] Two banks of $256\times16$ bits. The datapath reads
  the \emph{active} bank and the bus writes the \emph{shadow} bank. On a commit
  request, the banks swap at the next frame boundary (\cref{app:registers}).
  \item[\texttt{tilt\_gain}] A running product: at $k=0$ load
  $g_\text{out}$, then multiply by $r$ once per bin (O2).
  \item[\texttt{min\_tracker}] Combinational per-bin logic around the stored
  minima (\cref{lst:rtl:mintracker}). The sub-window frame counter and slot
  pointer are global, so they live in \code{frame\_ctrl}, not per bin.
\end{description}

% =========================================================================
\section{RTL coding guidelines}\label{sec:rtl:style}

Pick these rules on day one and never break them:
\begin{itemize}
  \item Synthesizable SystemVerilog only~\cite{sutherland2006svdesign}:
        \code{always\_ff} for registers, \code{always\_comb} for logic, and
        non-blocking assignments (\code{<=}) only in \code{always\_ff}
        blocks~\cite{cummings2000nonblocking}.
  \item One reset style: synchronous, active-high \sig{rst}.
  \item Naming: registers end in \sig{_q} with next state \sig{_d}; widths are
        parameters, never magic numbers.
  \item No latches: every \code{always\_comb} assigns every output on every
        path. Assign defaults at the top.
  \item Lint clean with \code{verilator --lint-only -Wall} before every commit.
        A warning you do not understand is a bug you have not found yet.
  \item Rounding and saturation are written once, in \code{qmul\_sat} and
        similar helpers, and match \file{golden/fixedpoint.py} exactly.
\end{itemize}

\begin{lstlisting}[style=sv,caption={\texttt{qmul\_sat}: the hardware twin of \texttt{fixedpoint.mul}.},label={lst:rtl:qmul}]
// Saturating fixed-point multiply with round-half-up, bit-exact with
// golden/fixedpoint.py:mul.  Assumes SHIFT = A_F + B_F - O_F >= 1.
module qmul_sat #(
  parameter int A_W = 16, A_F = 15,
  parameter int B_W = 16, B_F = 13,
  parameter int O_W = 16, O_F = 13
) (
  input  logic signed [A_W-1:0] a,
  input  logic signed [B_W-1:0] b,
  output logic signed [O_W-1:0] y
);
  localparam int P_W   = A_W + B_W;
  localparam int SHIFT = A_F + B_F - O_F;
  localparam logic signed [P_W-1:0] HALF = 1 <<< (SHIFT - 1);
  localparam logic signed [P_W-1:0] MAXV = (1 <<< (O_W - 1)) - 1;
  localparam logic signed [P_W-1:0] MINV = -(1 <<< (O_W - 1));

  logic signed [P_W-1:0] p, r;
  always_comb begin
    p = a * b;                    // exact product
    r = (p + HALF) >>> SHIFT;     // round half up, arithmetic shift
    if      (r > MAXV) y = MAXV[O_W-1:0];
    else if (r < MINV) y = MINV[O_W-1:0];
    else               y = r[O_W-1:0];
  end
endmodule
\end{lstlisting}

\begin{lstlisting}[style=sv,caption={\texttt{min\_tracker}: per-bin logic of the two-stage minimum tracker.},label={lst:rtl:mintracker}]
// One bin per cycle. Log-domain values are stored with a bias so they are
// all non-negative: unsigned comparison is then exact (and cheap).
module min_tracker #(parameter int W = 8, parameter int U = 8) (
  input  logic                 sub_end,    // last frame of a V-frame sub-window
  input  logic [$clog2(U)-1:0] slot,       // sub-window slot to overwrite
  input  logic [W-1:0]         x,          // this frame's smoothed value
  input  logic [W-1:0]         running_q,  // stored running minimum
  input  logic [U-1:0][W-1:0]  store_q,    // stored sub-window minima
  output logic [W-1:0]         running_d,
  output logic [U-1:0][W-1:0]  store_d,
  output logic [W-1:0]         s_min       // = Python's np.minimum(store.min(), running)
);
  logic [W-1:0] run_new, m;
  always_comb begin
    run_new   = (x < running_q) ? x : running_q;
    store_d   = store_q;
    running_d = run_new;
    if (sub_end) begin
      store_d[slot] = run_new;    // close the sub-window
      running_d     = x;          // restart from the current value
    end
    m = running_d;
    for (int i = 0; i < U; i++)
      if (store_d[i] < m) m = store_d[i];
    s_min = m;
  end
endmodule
\end{lstlisting}

\begin{tipbox}[Read your RTL as hardware]
The \code{for} loop in \cref{lst:rtl:mintracker} is not a loop in time. It
describes eight comparators and multiplexers, which synthesis builds as a chain.
At a few MHz a chain is fine. If timing ever fails, rewrite it as a balanced tree
of depth $\log_2 9\approx4$. Always ask \enquote{what circuit does this
describe?}~\cite{harris2021ddca}.
\end{tipbox}

% =========================================================================
\section{Build order}\label{sec:rtl:steps}

Build from the leaves up. Each block gets a cocotb test, comparing it with the
golden model, \emph{before} the next block is started
(\cref{ch:verification}).

\begin{stepbox}[label=st:rtl:skeleton]{Skeleton, lint and synthesis flow}
Create \file{rtl/} with empty modules and the port lists of
\cref{fig:rtl:arch}, a \code{Makefile} target for Verilator lint, and a Yosys
script that synthesizes any module and prints \code{stat}. Add both to CI.

\textbf{Check:} \code{make lint} and \code{make synth TOP=ise\_top} run cleanly on
the skeleton. \textbf{Effort:} half a day.
\end{stepbox}

\begin{stepbox}[label=st:rtl:qmul]{\texttt{qmul\_sat} and the other arithmetic helpers}
Implement \cref{lst:rtl:qmul} and the matching add/saturate helpers.

\textbf{Check:} an exhaustive or large random cocotb test against
\code{fixedpoint.mul} gives zero mismatches. Include the corner cases: both
operands at their most negative value, and exact half-LSB ties.
\textbf{Effort:} 1 day, including your first real cocotb test.
\end{stepbox}

\begin{stepbox}[label=st:rtl:recip]{\texttt{recip\_unit} and \texttt{exp\_unit}}
Implement the leading-zero normaliser, seed ROM, Newton--Raphson step and
pipeline registers. Generate the ROM contents from Python with the same function
the golden model uses, and write them as \code{.hex} files read with
\code{\$readmemh}.

\textbf{Check:} bit-exact over the full input range, or a dense random sample
of it, for both units. \textbf{Effort:} 2--3 days.
\end{stepbox}

\begin{stepbox}[label=st:rtl:lutregs]{\texttt{gain\_lut}, \texttt{tilt\_gain} and \texttt{regs\_axil}}
Implement the register map of \cref{app:registers}, the double-buffered table
and the frame-boundary commit. A minimal AXI4-Lite slave has one write channel
pair and one read channel. Handle the write address and data arriving in either
order.

\textbf{Check:} a cocotb test writes a table and commits it, and the new values
appear at the next frame boundary and not before. A second commit without
rewriting the table is rejected or handled exactly as specified.
\textbf{Effort:} 3--4 days.
\end{stepbox}

\begin{stepbox}[label=st:rtl:mintrk]{\texttt{min\_tracker} and \texttt{state\_ram}}
Implement \cref{lst:rtl:mintracker}. Infer \code{state\_ram} as a simple
synchronous RAM: write the RTL so synthesis recognises a block RAM, and check
the synthesis report to confirm.

\textbf{Check:} over many frames of golden vectors, $S_\text{min}$ and
$\tilde S_\text{min}$ match the model exactly, including across sub-window
boundaries and after pre-roll seeding. \textbf{Effort:} 2 days.
\end{stepbox}

\begin{stepbox}[label=st:rtl:binpipe]{\texttt{bin\_pipeline}, fed with golden spectra}
Assemble the per-bin pipeline (\cref{tab:rtl:stages}). Do not wait for the FFT:
drive the pipeline directly with the golden model's spectra and compare its
outputs and internal probes frame by frame.

\textbf{Check:} bit-exact for at least one full utterance and the whole
synthetic set, including a parameter change. Record the pipeline latency and the
cycles per frame. \textbf{Effort:} 1.5--2 weeks. This is the biggest block.
\end{stepbox}

\begin{stepbox}[label=st:rtl:fft]{\texttt{fft\_r2\_mem}}
Implement the memory-based FFT: complex first, block floating point included,
inverse by conjugation. Test it alone against \file{golden/fx_stft.py} with
impulses (every output bin is equal), single tones (one bin), and random input.

\textbf{Check:} bit-exact forward and inverse transforms; a forward transform
followed by the inverse returns the input within the documented rounding error.
\textbf{Effort:} 1--1.5 weeks.
\end{stepbox}

\begin{stepbox}[label=st:rtl:top]{\texttt{stft\_frontend}, \texttt{ola\_backend}, \texttt{frame\_ctrl}, \texttt{ise\_top}}
Add the ping-pong input buffer and window ROM, the synthesis window and
overlap-add, and the frame scheduler with its per-phase cycle counters. Connect
everything in \code{ise\_top}.

\textbf{Check:} end-to-end bit-exact on the golden vectors and on a full
VoiceBank utterance. The pre-roll behaves as specified, and the cycle counters
report the frame schedule. \textbf{Effort:} 1 week.
\end{stepbox}

\begin{stepbox}[label=st:rtl:resources]{Resource log}
Synthesize \code{ise\_top} with Yosys (\code{stat}) and with Vivado
out-of-context synthesis for the Zynq-7020. Record LUTs, flip-flops, DSP slices
and block RAMs per block in \file{docs/resources.md}, with the commit hash.

\textbf{Check:} the design fits comfortably in the Zynq-7020, and block RAM use
matches your state-memory calculation. If it does not, find out why: memories
inferred as flip-flops are a common surprise. \textbf{Effort:} half a day,
then 15 minutes per milestone.
\end{stepbox}

% =========================================================================
\section{Pitfalls}\label{sec:rtl:pitfalls}

\begin{pitfallbox}[Pitfalls in Phase 2]
\begin{itemize}
  \item \textbf{Writing the whole design before testing anything.} Bugs
        multiply when blocks are combined. Test every block alone first.
  \item \textbf{Signed/unsigned mix-ups.} In SystemVerilog, one unsigned operand
        makes the whole expression unsigned. Slices of signed vectors are
        unsigned. Be explicit with \code{\$signed()} and check with lint.
  \item \textbf{Off-by-one at the spectrum edges.} The 5-tap smoothing at bins 0,
        1, $K-2$ and $K-1$ must zero-pad exactly like \code{np.convolve}.
  \item \textbf{Parameter updates mid-frame.} Any register the datapath reads
        must be shadowed and latched at the frame boundary. Otherwise one frame
        mixes old and new parameters, and you will hear it as a click.
  \item \textbf{Inferred latches and flip-flop memories.} Read the synthesis
        report after every major change.
  \item \textbf{Optimising early.} Get it bit-exact first, then make it smaller.
\end{itemize}
\end{pitfallbox}

\begin{readbox}
Digital design fundamentals: Harris and Harris~\cite{harris2021ddca}.
Architecture transformations for DSP (pipelining, folding, retiming):
Parhi~\cite{parhi1999vlsidsp} and Kaeslin~\cite{ch12-13:kaeslin2014topdown}.
FPGA arithmetic: Meyer-Baese~\cite{meyerbaese2014fpgadsp}. FFT hardware
architectures: Garrido's survey~\cite{ch12-13:garrido2022fft}. Practical
Verilog and formal-friendly design: the ZipCPU blog~\cite{ch12-13:zipcpu}.
\end{readbox}

\begin{donebox}[Phase 2 is done (with Phase 3) when]
\begin{checklist}
  \item Every module in \cref{fig:rtl:arch} exists, is lint-clean, and has its own passing test.
  \item \code{ise\_top} is bit-exact with the golden model on all golden vectors and at least one full utterance.
  \item Cycle counters report the frame schedule; H4 can be evaluated from measured cycles.
  \item The resource log is up to date and the design fits the Zynq-7020.
\end{checklist}
\end{donebox}
